\chapter{Introduction}\label{ch:intro}

\IS is a 64-bit load/store architecture with fixed 32-bit instructions, built on one idea:
\emph{memory beyond the core is explicit I/O, and every storage tier is visible to software.}

\section{Design principles}

\begin{description}[leftmargin=1em,style=nextline]
\item[Loads and stores reach only the Local Store.] Each core has a private data Local Store
  (D-LS), the only memory that load and store instructions can address, and a private instruction
  Local Store (I-LS), the only memory from which instructions are fetched: 4 to 64~KiB each, 64~KiB by
  default. The Local Stores are architectural scratchpads at the latency of a first-level cache.
\item[Everything else is reached by transfers.] The cluster store (T2), system memory (T3) and devices
  (T4) are reached only through the core's \term{Transfer Unit} (TU), by block copies
  (\insn{tget}, \insn{tput}, \insn{tgetx}) or by far atomics (\insn{tamo}). There are no hardware
  caches and nothing to keep coherent; software builds caches when it wants them, with hardware assists
  for the tag lookup (\insn{ltag}, \insn{linst}) and for code overlays (\insn{xmap}, \insn{xcall}).
\item[Far memory is named by capabilities.] Every transfer and atomic names its target segment through
  a capability register: a hardware-protected (space, base, length, permissions) descriptor that
  software can split, join, restrict, inspect and pass between cores, but never forge.
\item[Writable capabilities are linear.] A capability that permits writing cannot be copied; moving it
  empties its source, and no instruction may overwrite one that is still live. At any instant a byte of
  memory has a single writer, or any number of readers, or any number of atomic users. Data races between
  cores cannot be expressed.
\item[Transfers are as-if-synchronous.] Each transfer takes effect at its issue point in program order.
  The hardware stalls any later access to a Local Store range that is still in flight. Overlapping
  transfers with computation is a performance decision, never a correctness one, and a single core's
  results do not depend on timing.
\item[Ownership transfer is synchronisation.] Cores communicate by giving capabilities to each other's
  inboxes (\insn{cgive}, \insn{ctake}). A gift completes only after the giver's memory transfers have
  finished, so handing over a buffer also publishes its contents.
\item[Authority can be lent and taken back.] A region of memory can be made revocable
  (\insn{cdom}). Everything derived from it, however it was split, shared or passed on, dies at once when
  the holder of its revoker revokes it (\insn{crevoke}, \insn{cundom}); revocation works by address, on
  256~KiB granules, so the number of revocable regions is limited only by memory.
\item[Two privilege modes, no interrupts.] A supervisor runs programs in user mode, where the
  capability save table, the revocation operations and the core window are out of reach. There are no
  interrupts: a device signals an event by completing a transfer, a core waits for any of several events
  with \insn{tany}, and the only forced preemption is a stop request sent through core control. A program
  asks its supervisor for service with \insn{svc}.
\item[Time and pictures are devices like any other.] The timer is a set of device registers, and waiting
  for a deadline is a read that completes when the time has come. The framebuffer is ordinary system
  memory that the display scans out, owned and passed around through capabilities like any other
  memory.
\end{description}

\section{System organisation}

Figure~\ref{fig:system} shows an \IS system. Each core contains its integer and capability registers,
its two Local Stores and its Transfer Unit. All cores share the far tiers through an arbiter. The block
device is itself a node that can receive and return capabilities.

\begin{figure}[htbp]
\centering
\begin{tikzpicture}[
  box/.style={draw=accent,fill=accentlight,rounded corners=2pt,minimum height=8mm,align=center,font=\small},
  ls/.style={draw=accent,fill=white,minimum width=19mm,minimum height=7mm,align=center,font=\footnotesize},
  far/.style={draw=black!60,fill=black!5,minimum height=9mm,align=center,font=\small},
  >={Stealth[length=2.2mm]}]
  \foreach \i/\x in {0/0,1/4.6,2/9.2} {
    \node[draw=accent,line width=0.8pt,minimum width=40mm,minimum height=38mm,rounded corners=3pt] (core\i) at (\x,0) {};
    \node[font=\small\bfseries,anchor=north] at (core\i.north) {\ifnum\i=2 Core $N-1$\else Core \i\fi};
    \node[box,minimum width=34mm] (rf\i) at (\x,0.75) {x0--x31 \quad c0--c15};
    \node[ls] (ils\i) at (\x-0.95,-0.35) {I-LS 64 KiB};
    \node[ls] (dls\i) at (\x+0.95,-0.35) {D-LS 64 KiB};
    \node[box,minimum width=34mm,minimum height=6mm] (tu\i) at (\x,-1.35) {Transfer Unit, tg0--tg15};
    \draw[<->] (rf\i.south -| dls\i) -- node[right,font=\scriptsize]{ld/st} (dls\i.north);
    \draw[->] (ils\i.north) -- node[left,font=\scriptsize]{fetch} (rf\i.south -| ils\i);
    \draw[<->] (tu\i.north -| dls\i) -- (dls\i.south);
    \draw[->] (tu\i.north -| ils\i) -- (ils\i.south);
  }
  \node at (6.9,0) {\Large$\cdots$};
  \node[draw=black!60,minimum width=133mm,minimum height=7mm,font=\small] (arb) at (4.6,-3.0) {Far arbiter (TU transfers and far atomics)};
  \foreach \i in {0,1,2} \draw[<->] (tu\i.south) -- (tu\i.south |- arb.north);
  \node[far,minimum width=30mm] (t2) at (0,-4.6) {T2 cluster store\\256 KiB};
  \node[far,minimum width=30mm] (t3) at (4.6,-4.6) {T3 system memory\\64 MiB};
  \node[far,minimum width=30mm] (t4) at (9.2,-4.6) {T4 device space\\console, disk, timer, display};
  \foreach \t in {t2,t3,t4} \draw[<->] (\t.north) -- (\t.north |- arb.south);
  \draw[dashed,accent,->,thick] (core0.north) to[bend left=25] node[above,font=\scriptsize]{cgive / ctake (inboxes)} (core1.north);
\end{tikzpicture}
\caption{An \IS system. Solid arrows carry data; the dashed arrow carries capabilities.}
\label{fig:system}
\end{figure}

\section{Storage tiers}

\begin{table}[htbp]
\centering\small
\caption{Storage tiers of the prototype platform}\label{tab:tiers}
\begin{tabularx}{\linewidth}{@{}l l l X@{}}
\toprule
\textbf{Tier} & \textbf{Name} & \textbf{Size} & \textbf{Access}\\\midrule
T1i & Instruction Local Store (I-LS) & 64 KiB per core (4--64 KiB) & Instruction fetch; written only by \insn{tgetx}\\
T1d & Data Local Store (D-LS) & 64 KiB per core (4--64 KiB) & Loads and stores; \insn{tget}/\insn{tput} endpoints\\
T2 & Cluster store & 256 KiB, shared & Transfers and far atomics through a capability (space 2)\\
T3 & System memory & 64 MiB, shared & Transfers and far atomics through a capability (space 3); the top 16 MiB is the framebuffer window\\
T4 & Device space & \hex{5000} bytes & Transfers of one byte per beat through a device capability (space 4)\\
\bottomrule
\end{tabularx}
\end{table}

Two further spaces hold no storage: core-control space (space 5), through which a core stops other
cores, and endpoint space (space 6), whose capabilities name nodes.

The sizes are those of the prototype platform in its default configuration. Every size is a property
of the platform, reported by the platform table (Chapter~\ref{ch:platform}). The Local Store sizes,
the \texttt{DLS\_SIZE} and \texttt{ILS\_SIZE} of the pseudocode, are powers of two from 4~KiB to
64~KiB, chosen independently (Section~\ref{sec:lssize}); every core of a machine has the same ones.

\section{A first program}

The smallest complete program prints a greeting from every core in turn. It shows the three things
that make \IS different: code and data are in the Local Stores, output is a transfer through a
capability, and cores are sequenced by passing a capability.

\begin{asm}
        .data                           ; data Local Store
msg:    .ascii "Hello from Strata core "
msg_end:
        .text                           ; instruction Local Store; execution starts at 0
_start: csrr  s0, coreid                ; s0 = this core's number
        beqz  s0, speak                 ; core 0 speaks first
        addi  t0, s0, -1                ; t0 = the core before us
        ctake c3, t0                    ; wait for its token (only from that core)
speak:  la    a0, msg                   ; D-LS address of the text
        li    a1, msg_end - msg         ; its length
        tput  tg15, a0, c2, zero, a1    ; console output is a transfer through c2
        ...                             ; (print the core number, a newline)
        addi  t0, s0, 1                 ; t0 = the next core (the last core skips this)
        cgive c2, t0                    ; pass a copy of the console capability on;
                                        ; completes only after our output has gone out
        halt
\end{asm}

Chapter~\ref{ch:examples} walks through this and the other example programs in detail.

\section{Terminology}
\begin{description}[leftmargin=1em,style=nextline]
\item[Node] A destination of \insn{cgive}: cores $0 \ldots N-1$ and block device 0 as node $N$, where
  $N$ is the number of cores.
\item[Segment] The contiguous range of a far space named by a capability.
\item[In flight] A transfer that has issued but whose data movement has not finished. In-flight state
  is never architecturally observable except through timing-visible state (Section~\ref{sec:timing-visible}).
\item[Drain] To wait until all of the core's in-flight transfers have completed.
\item[Supervisor, user] The two privilege modes. Cores start in supervisor mode; programs run by an
  operating system run in user mode (Chapter~\ref{ch:traps}).
\item[Pending device read] A \insn{tget} from a device port that has nothing to deliver yet. It stays in
  flight until the device completes it; this is how a core waits for input, for Ctrl-C, for a deadline
  or for the display (Section~\ref{sec:events}).
\item[Revocation by address] A revocable capability carries the generation of the granule its base
  lies in; revoking the granule makes all of them unusable (Section~\ref{sec:domains}).
\item[Endpoint] A capability that names a node, used to give and take capabilities without naming the
  node by number (Section~\ref{sec:endpoints}).
\end{description}
A full glossary is in Appendix~\ref{app:glossary}.
